\subsection*{43. A Complete Prepaid-Cohort Example: Assumptions}

Consider a new, hypothetical entity conducting 1 3-month course. On January 1, the founder contributes \textbf{₹5,00,000 as share capital}. 10 learners collectively prepay \textbf{₹3,00,000} for the 3-month program, with no refunds, tax, or collection complications assumed. Assume the service obligation is satisfied evenly over January, February, and March, allowing \textbf{₹1,00,000 recognized revenue each month} for this teaching exercise.

The entity pays ₹1,20,000 of attributable acquisition and marketing expense in January, ₹60,000 for equipment with an assumed 12-month useful life and straight-line depreciation of ₹5,000 per month starting January, and ₹60,000 of monthly delivery expense (instructors, support, and platform) for each of the 3 months. Assume all expenditures and teaching occur as scheduled, with no other assets or liabilities, no corporate income tax, no financing charges, and no residual value on equipment.

These artificial assumptions let us reconcile the statements exactly. They are not a recommendation on real revenue schedules, depreciation rates, tax policy, or startup spending.

\subsection*{44. January's Income Statement: A Loss Despite Cash in the Bank}

In January, the firm recognizes ₹1,00,000 of earned course revenue. It incurs ₹60,000 of teaching and platform expense, ₹1,20,000 of acquisition and marketing expense, and ₹5,000 depreciation. The resulting \textbf{January net loss is ₹85,000} before any tax complexities (Table~\ref{tab:january-income-statement}).

\begin{table}[htbp]
\centering
\small
\begin{tabularx}{\columnwidth}{Xr}
\toprule
\textbf{January Income Statement Line} & \textbf{Amount} \\
\midrule
Course revenue recognized & ₹1,00,000 \\
Teaching and platform delivery expense & $-$₹60,000 \\
Marketing and acquisition expense & $-$₹1,20,000 \\
Depreciation expense & $-$₹5,000 \\
\midrule
\textbf{Net accounting result (loss)} & \textbf{$-$₹85,000} \\
\bottomrule
\end{tabularx}
\caption{Illustrative January accrual income statement.}
\label{tab:january-income-statement}
\end{table}

Some would mistakenly report ₹3,00,000 as January revenue because that much tuition arrived. Others might expense the entire ₹60,000 equipment purchase and omit depreciation. Either method would misrepresent the simplified accrual performance assumed here.

\subsection*{45. January's Balance Sheet Must Balance}

The company began with ₹5,00,000 cash and received ₹3,00,000 advance tuition. It then paid ₹1,20,000 marketing, ₹60,000 equipment, and ₹60,000 January delivery, leaving \textbf{₹5,60,000 cash}. Equipment has a \textbf{₹55,000 net carrying amount} after January depreciation. The company still owes 2 months of instruction, represented by a simplified \textbf{₹2,00,000 contract liability}. Equity equals original contributed capital ₹5,00,000 minus January loss ₹85,000, or \textbf{₹4,15,000} (Table~\ref{tab:january-balance-sheet}).

\begin{table}[htbp]
\centering
\small
\begin{tabularx}{\columnwidth}{Xr}
\toprule
\textbf{Balance Sheet Position (January 31)} & \textbf{Amount} \\
\midrule
Cash and cash equivalents & ₹5,60,000 \\
Equipment, net of accumulated depreciation & ₹55,000 \\
\midrule
\textbf{Total assets} & \textbf{₹6,15,000} \\
\midrule
Contract liability: unearned tuition & ₹2,00,000 \\
Share capital plus retained earnings & ₹4,15,000 \\
\midrule
\textbf{Total liabilities plus equity} & \textbf{₹6,15,000} \\
\bottomrule
\end{tabularx}
\caption{Illustrative January 31 balance sheet position.}
\label{tab:january-balance-sheet}
\end{table}

The 2 sides agree. The large cash balance is not proof that the course has earned all its tuition. Some cash is economically committed to fulfilling the remaining teaching obligations.

\subsection*{46. January's Cash Flow Statement Tells a Different Story}

Cash flow from operations includes ₹3,00,000 prepaid customer cash minus ₹1,20,000 marketing and ₹60,000 delivery: \textbf{₹1,20,000 operating cash inflow}. Investing activities include ₹60,000 outflow for equipment. Financing activities include the founder's ₹5,00,000 equity contribution. Total January cash is therefore \textbf{₹5,60,000} (Table~\ref{tab:january-cash-flow}).

\begin{table}[htbp]
\centering
\small
\begin{tabularx}{\columnwidth}{Xr}
\toprule
\textbf{Cash Flow Statement (January)} & \textbf{Amount} \\
\midrule
Operating cash inflow: advance tuition & ₹3,00,000 \\
Operating cash outflow: marketing and acquisition & $-$₹1,20,000 \\
Operating cash outflow: delivery and platform & $-$₹60,000 \\
\midrule
\textbf{Net cash from operating activities} & \textbf{₹1,20,000} \\
Investing cash outflow: capital equipment & $-$₹60,000 \\
Financing cash inflow: contributed share capital & ₹5,00,000 \\
\midrule
\textbf{Net increase in cash and closing balance} & \textbf{₹5,60,000} \\
\bottomrule
\end{tabularx}
\caption{Illustrative January cash flow classification.}
\label{tab:january-cash-flow}
\end{table}

The same result can be reconciled from accounting profit in this simple case: January net loss ($-$₹85,000) plus noncash depreciation (₹5,000) plus increase in contract liability (₹2,00,000) = \textbf{₹1,20,000 operating cash flow}. The arithmetic explains why cash increased despite an accounting loss.

\subsection*{47. The Full 3-Month Story Reveals the Economics}

By the end of March, the program has delivered all 3 months. Total recognized revenue is ₹3,00,000. Marketing expense is ₹1,20,000; teaching and platform expense is ₹1,80,000; depreciation over 3 months is ₹15,000. Cumulative accounting result is therefore \textbf{a ₹15,000 loss}.

Cash equals original equity plus prepaid tuition minus marketing, equipment, and 3 months of delivery:
\begin{equation}
\begin{aligned}
\text{Cash} &= ₹5,00,000 + ₹3,00,000 - ₹1,20,000 \\
&\quad - ₹60,000 - ₹1,80,000 = \textbf{₹4,40,000}.
\end{aligned}
\end{equation}
Equipment net book value is ₹45,000; the contract liability is zero; equity is ₹4,85,000. The balance sheet again balances at ₹4,85,000.

This 3-month course generates \textbf{zero cash contribution before equipment investment and depreciation} from the customer transactions: ₹3,00,000 tuition less ₹1,20,000 acquisition and ₹1,80,000 delivery. The founder's ₹5 lakh financing masks that weakness if 1 looks only at the bank account. The case demonstrates why funding a venture, collecting prepaid money, earning revenue, and operating profitably are 4 separate events.
